%Document-Author: Cognome Nome + Cognome Nome + ...
%Document-Date: AAAA-MM-GG
%Document-Description: Descrizione documento


\documentclass[a4paper]{report}
\usepackage[english, italian]{babel}
\usepackage[T1]{fontenc}
\usepackage[utf8]{inputenc}
\usepackage{url}
\usepackage{graphicx}
\graphicspath{{../Figures/}}
\usepackage[hidelinks]{hyperref}
\usepackage{booktabs}
\usepackage{tabularx}
\usepackage{pifont}
\usepackage[table]{xcolor}
\usepackage{float}

\newcolumntype{s}{>{\hsize=.21\hsize}X}
\newcolumntype{f}{>{\hsize=.37\hsize}X}
\newcolumntype{m}{>{\hsize=.42\hsize}X}

\newcommand{\mychapter}[2]{
	\setcounter{chapter}{#1}
	\setcounter{section}{0}
	\setcounter{subsection}{1}
	\chapter*{#2}
	\addcontentsline{toc}{chapter}{#2}
}

\renewcommand{\abstractname}{Tabella contenuti}

\begin{document}
	
	\begin{titlepage}
		% Defines a new command for the horizontal lines, change thickness here
		\newcommand{\HRule}{\rule{\linewidth}{0.5mm}} 
		\center  
		
		% HEADING SECTION
		\textsc{\LARGE KaleidosCode}\\[1.5cm] 
		\textsc{\Large SweDesigner}\\[0.5cm] 
		\textsc{\large Software per diagrammi UML}\\[0.5cm]
		
		
		% TITLE SECTION
		\HRule \\[0.4cm]
		{ \huge \bfseries Norme di progetto}\\[0.4cm] 
		\HRule \\[1.5cm]
		
		% AUTHOR SECTION
		\begin{minipage}{0.4\textwidth}
			\begin{flushleft} \large
				\emph{Redattori:}\\
				Bonato Enrico\\
				Bonolo Marco\\
				Pace Giulio\\
				Pezzuto Francesco\\
			\end{flushleft}
		\end{minipage}
		~
		\begin{minipage}{0.4\textwidth}
			\begin{flushright} \large
				\emph{Approvazione:} \\
				Sovilla Matteo\\
				\emph{Verifica:} \\
				Sanna Giovanni
			\end{flushright}
		\end{minipage}
		
		%immagine
		\begin{figure}[H]
			\centering
			\includegraphics[width=\textwidth]{Figures/KaleidosCodeLogo.png}
		\end{figure}
		
		%destinazione d'uso e lista di distribuzione	
		\begin{minipage}{0.4\textwidth}
			\begin{flushleft} \large
				\emph{Uso:}\\
				Interno \\
				
			\end{flushleft}
		\end{minipage}
		~
		\begin{minipage}{0.4\textwidth}
			\begin{flushright} \large
				\emph{Lista di distribuzione:} \\
				Prof. Vardanega Tullio\\
				Prof. Cardin Riccardo\\
				Zucchetti spa %Controllare sia scritto giusto
			\end{flushright}
		\end{minipage}
		
		\begin{center}
			Versione 1.0.0
		\end{center}
		% Date, change the \today to a set date if you want to be precise
		{\large \today}\\[3cm] 
		% Fill the rest of the page with whitespace
		\vfill  
	\end{titlepage}
	
	
	\tableofcontents
	
	\mychapter{0}{Diario delle modifiche}
		\begin{table}[H]
			\begin{tabularx}{\textwidth}{s f m X}
				\noalign{\hrule height 1.5pt}
				\rowcolor{orange!85} Versione & Data & Autore & Descrizione \\
				\noalign{\hrule height 1.5pt}
				1.0.0 & 2017/03/01 & Pace Giulio & Creazione scheletro del documento e stesura parziale dei documenti\\
			\end{tabularx}
			\caption{Diario delle modifiche \label{tab:table_label}}
		\end{table}

	\mychapter{1}{Introduzione}
	\section{Scopo del documento}
		\dots
	\section{Scopo del prodotto}
		\dots
		\section{Glossario}
		\dots
	\section{Riferimenti}
	\subsection{Informativi}	
		\dots
	
	\mychapter{2}{Comunicazioni}
	
	\mychapter{3}{Riunioni}
	\section{Frequenza}
		Le riunioni interne del gruppo di lavoro avranno una cadenza settimanale.
		\dots
	\section{Convocazione riunione}
		\subsection{Interna}
			Il Responsabile di Progetto ha il compito di convocare le riunioni generali interne, ossia
			le riunioni a cui sono convocati tutti i membri del gruppo.
			Se un componente qualsiasi ritiene necessaria la convocazione di una riunione, deve
			inoltrare la richiesta al responsabile il quale decide se respingere o accettare
			tale richiesta.
			Il responsabile deve convocare l'assemblea con almeno quattro giorni di preavviso
			attraverso l'invio di una richiesta nella sezione Comunicazioni di Asana, il cui
			corpo è formato da:
			\begin{itemize}
				\item Oggetto: Convocazione riunione n. X (dove X indica il numero progressivo
				di riunioni effettuate)
				\item Corpo:
				\begin{itemize}
					\item Data e ora prevista
					\item Luogo previsto
					\item Ordine del giorno
				\end{itemize}
			\end{itemize}
			Ogni componente deve rispondere al messaggio nel minor tempo possibile confermando
			la presenza o in caso contrario motivando l'eventuale assenza. In caso di manca risposta
			il Responsabile di Progetto è obbligato a contattare direttamente il componente che non
			ha fornito una risposta. Una volta ricevute tutte le risposte, il Resposanbile di Progetto
			può decidere se confermare, rinviare o annullare la riunione, basandosi sul numero di
			presenze e/o assenze, per garantire il massimo numero di presenti. La conferma, il rinvio,
			l'annullamento devono essere effettuati tramite messaggio nella sezione Comunicazioni di cui
			sopra.
			\dots
		\subsection{Esterna}
			Il Responsabile di Progetto ha il compito di convocare le riunioni generali esterne, ossia
			le riunioni a cui sono convocati tutti i membri del gruppo e il proponente e/o committente.
			Se un componente qualsiasi ritiene necessaria la convocazione di una riunione, deve
			inoltrare la richiesta al responsabile il quale decide se respingere o accettare
			tale richiesta.
			Il responsabile deve prima di tutto concordare con il proponente e/o committente delle date e dei luoghi
			possibili in cui svolgere la riunione attraverso l'invio di un email contenente:
			\begin{itemize}
				\item Oggetto: Richiesta riunione
				\item Corpo
				\begin{itemize}
					\item Motivazione
					\item Eventuali date e/o luoghi possibili
				\end{itemize}
			\end{itemize}
			Dopo aver concordato delle date possibili, il responsabile deve inviare un messaggio a
			i membri del gruppo nella sezione Comunicazioni in Asana in cui sono specificate:
			\begin{itemize}
				\item Oggetto: Richiesta riunione col proponente
				\item Corpo
				\begin{itemize}
					\item Date e/o luoghi possibili
				\end{itemize}
			\end{itemize}
			Ogni membro del gruppo deve rispondere alla comunicazione confermando la presenza o in
			caso contrario motivando l'eventuale essenza. In caso di mancata risposta, il Responsabile
			di Progetto deve contattare direttamente il componente che non ha fornito una risposta.
			Una volta ricevute tutte le risposte, il responsabile può decidere se confermare, rinviare,
			annullare la riunione con il proponente basandosi sul numero di presenze e /o assenze.
			Il Responsabile di Progetto deve poi, in caso di:
			\begin{itemize}
				\item conferma: comunicare a tutti i membri e al proponente e/o committente la data e il luogo definitivi
				\item rinvio: comunicare a tutti i membri la decisione presa e ripetere il procedimento
				dall'inizio
				\item annullamento: comunicare a tutti i membri la decisione presa e ripetere il
				procedimento dall'inizio
			\end{itemize}
	\section{Svolgimento riunione}
		\subsection{Esterna}
			All'inizio di ogni riunione, verificata la presenza dei membri previsti, viene scelto un
			segretario che avrà il compito di annotare gli argomenti trattati e di redigere il verbale
			della riunione, che dovrà poi essere inviato ai restanti membri del gruppo.
			Tutti i partecipanti devono tenere un comportamento consono al miglior svolgimento
			dell'assemblea e al raggiungimento degli obbiettivi della stessa. Il segretario deve inoltre assicurarsi che venga seguito l'ordine del giorno in modo da affrontare ogni argomento previsto.
			\dots
			\section{Verbale}
			\subsection{Riunione interna}
			Il verbale di Riunione interna è un documento informale che permette di tenere traccia
			degli argomenti discussi in ogni riunione.
			Il segretario, scelto a rotazione, ha il compito di redigere questo documento.
			Il verbale prodotto deve poi essere condiviso con tutti i membri del gruppo con Google
			Drive, in modo da rendere sempre disponibile e accessibile il contenuto dello stesso.
			\dots
		\subsection{Riunione esterna}
			In caso di riunione esterna con il proponente e/o committente, il verbale è un documento
			ufficiale che può avere valore normativo e quindi deve essere redatto seguendo criteri
			specifici.
			Per agevolare la scrittura di tale documento è stato preparato un template LATEX che
			ne definisce la struttura e ne organizza i contenuti. Vi è quindi l'obbligo di seguire
			il sopraccitato schema per creare il verbale, di inviare il verbale prodotto come allegato al 
			proponente e/o committente in risposta all'email della riunione e di condividere il 
			verbale prodotto con tutti i membri del gruppo con Google Drive.
			\dots
	
	\mychapter{4}{Documenti}
	
	\mychapter{5}{Analisi dei requisiti}
	
	\mychapter{6}{Codifica dei file e documentazione}
	
	\mychapter{7}{Glossario}
	
	\mychapter{8}{Protocollo per lo sviluppo dell'applicazione}
	
	\mychapter{9}{Ambiente di lavoro}
		
	\cleardoublepage
	\addcontentsline{toc}{chapter}{\listfigurename}
	\listoffigures
	
	\cleardoublepage
	\addcontentsline{toc}{chapter}{\listtablename}
	\listoftables
		
\end{document}
	